%% kapa: exam=2 round=37 year=2026 period=2 subject=appraisal_theory subject_ko=감정평가이론 seq=4 number=4 points=10
%% source: Q-Net 감정평가사 기출문제 (cst00309) artlSeq=5268128 file=raw/1008/5268128/26년 제37회 감정평가사 2차 (2교시_감정평가이론).pdf
%% conversion: kordoc 4.24.51 -> tools/gen2.py (auto); review=proofread-gpt-terra; flags=none
\begin{longquestion}{【문제 4】}
감정평가에 있어 물적 불일치의 개념을 쓰고, 물적 불일치가 발생하는 토지 또는 건물의 사례를 3개 제시한 후, 그 처리방법을 쓰시오. (10점) \par
\end{longquestion}
